%% interpretability.tex  —  Moufida Interpretability Framework
%% pdflatex interpretability.tex

\documentclass[11pt,a4paper]{article}

\usepackage[margin=2.2cm]{geometry}
\usepackage{amsmath,amssymb}
\usepackage{graphicx}
\usepackage[hidelinks]{hyperref}
\usepackage{xcolor}
\usepackage{booktabs}
\usepackage{enumitem}
\usepackage{tcolorbox}
\usepackage{microtype}
\usepackage{titlesec}
\tcbuselibrary{skins,breakable}

%% ── Palette (section titles and box borders/backgrounds only) ───────────────
\definecolor{IntBlue}{RGB}{25,84,163}
\definecolor{IntTeal}{RGB}{0,115,100}
\definecolor{IntViolet}{RGB}{90,50,165}
\definecolor{IntSlate}{RGB}{50,62,75}
\definecolor{IntAmber}{RGB}{175,100,0}
\definecolor{BlueBg}{RGB}{240,246,255}
\definecolor{TealBg}{RGB}{238,248,245}
\definecolor{VioletBg}{RGB}{245,241,255}
\definecolor{AmberBg}{RGB}{255,248,235}
\definecolor{BoxBg}{RGB}{250,251,252}

%% ── Spacing ─────────────────────────────────────────────────────────────────
\setlength{\parskip}{4pt plus 1pt}
\setlength{\parindent}{0pt}
\setlist{itemsep=2pt, topsep=3pt, parsep=0pt, leftmargin=1.4em}
\setlist[enumerate]{itemsep=2pt, topsep=3pt}

%% ── Section styling ─────────────────────────────────────────────────────────
\titlespacing*{\section}{0pt}{14pt}{5pt}
\titlespacing*{\subsection}{0pt}{10pt}{3pt}
\titleformat{\section}{\large\bfseries\color{IntBlue}}{{\thesection}}{0.6em}{}[\vspace{-2pt}\color{IntBlue}\titlerule]
\titleformat{\subsection}{\normalsize\bfseries\color{IntSlate}}{{\thesubsection}}{0.6em}{}

%% ── Colour boxes ─────────────────────────────────────────────────────────────
\tcbset{
  base/.style={
    boxrule=0.6pt, arc=3pt,
    boxsep=4pt, top=4pt, bottom=4pt, left=6pt, right=6pt,
    coltitle=white, fonttitle=\small\bfseries, breakable
  }
}
\newtcolorbox{conceptbox}[1]{base,
  colback=BlueBg, colframe=IntBlue!75,
  title={#1}}
\newtcolorbox{mathbox}[1]{base,
  colback=TealBg, colframe=IntTeal!80,
  title={#1}}
\newtcolorbox{implbox}[1]{base,
  colback=VioletBg, colframe=IntViolet!70,
  title={#1}}
\newtcolorbox{notebox}{base,
  colback=AmberBg, colframe=IntAmber!80,
  coltitle=white, title={Note}}

%% ── Title ───────────────────────────────────────────────────────────────────
\title{%
  {\color{IntBlue}\rule{\linewidth}{1.5pt}}\\[6pt]
  \textbf{\Large Interpretability in Moufida}\\[4pt]
  {\normalsize Three Interlocking Layers of Transparent AI Reasoning}\\[6pt]
  {\color{IntBlue}\rule{\linewidth}{0.8pt}}
}
\author{Team Makrouna Kadheba}
\date{AINS 2026}

\begin{document}
\maketitle\thispagestyle{empty}

\begin{abstract}\small
Most AI scoring tools produce a number with no auditable chain of reasoning.
Moufida treats interpretability as a first-class system requirement: every axis score
must be decomposable, every retrieved evidence chunk must carry a relevance justification,
and every intake question must be selectable by a computable criterion.
Three interlocking techniques address these requirements at different pipeline layers —
\textbf{Concept Bottleneck Models} (why did this axis score $s$?),
\textbf{Contrastive Axis Direction Probes} (why was this evidence retrieved?),
and \textbf{Computerized Adaptive Testing / IRT} (why was this question asked?).
This document explains the academic foundation, the mathematics, and the concrete
implementation of each technique in Moufida.
\end{abstract}

\tableofcontents
\vspace{6pt}
\noindent\color{IntBlue}\rule{\linewidth}{0.4pt}

%% ─────────────────────────────────────────────────────────────────────────────
\section{The Interpretability Problem}
%% ─────────────────────────────────────────────────────────────────────────────

Startup scoring tools, whether LLM-based or regression-based, typically reduce a rich
profile to a single scalar. That number carries no causal structure: the founder cannot
trace which sub-dimension drove the outcome, why a particular document was cited,
or whether the questions asked were the right ones for their stage.

\begin{itemize}
  \item For high-stakes decisions — applying for a grant, hiring a CTO, pivoting the model —
        an unexplained recommendation is professionally insufficient.
  \item Tunisian founders who share Moufida's output with investors, incubators (APII,
        Smart Capital), or government bodies expect a traceable, auditable report.
  \item A score of ``58/100'' is useless without knowing: \emph{which concept} is suppressing
        it, \emph{which evidence} justifies the assessment, and \emph{how} the profile
        was understood in the first place.
\end{itemize}

Moufida's design principle is that every score must come with a \textbf{decomposable
audit trail}. This is enforced architecturally: three separate services each own one
interpretability layer.

\begin{center}\small
\begin{tabular}{@{}lll@{}}
\toprule
\textbf{Layer} & \textbf{Question answered} & \textbf{Service} \\
\midrule
Score decomposition & Why is axis $A$ scoring $s$? & \texttt{signal/} (Rust) \\
Evidence relevance  & Why was chunk $c$ retrieved for axis $A$? & \texttt{rag/} (Python) \\
Intake transparency & Why was question $j$ selected? & \texttt{orchestrator/} (Python) \\
\bottomrule
\end{tabular}
\end{center}

%% ─────────────────────────────────────────────────────────────────────────────
\section{Layer 1 — Concept Bottleneck Model}
%% ─────────────────────────────────────────────────────────────────────────────

\subsection{Academic Foundation}

\textbf{Koh et al.} — \textit{Concept Bottleneck Models.} ICML 2020.\\
\textbf{Oikarinen et al.} — \textit{Label-Free Concept Bottleneck Models.} ICLR 2023.

\begin{conceptbox}{The Core Idea}
A standard model maps input $x$ directly to output $y$ through opaque learned weights.
A \textbf{Concept Bottleneck Model} inserts a layer of human-named concepts
$\mathbf{c} = (c_1, \dots, c_k)$ as an obligatory intermediate: the model must first
score each concept, then combine those scores to produce $y$.
Because every concept has a human-readable name, you can inspect which one is holding
the score back and project the gain from improving it.
\end{conceptbox}

\begin{notebox}
Standard CBMs require a large labelled dataset of $(x_i, \mathbf{c}_i, y_i)$ triples.
Moufida has no historical corpus of Tunisian startups with ground-truth concept ratings.
\textbf{Label-Free CBM} (Oikarinen et al.) removes this requirement: an LLM acts as a
zero-shot concept scorer, replacing the annotation bottleneck.
\end{notebox}

\subsection{The Scoring Formula}

For each diagnostic axis, the orchestrator prompts the LLM to score 4–7 named concepts
for the current profile ($\hat{c}_i \in [0,1]$), then passes those scores to the
Affinitree scoring engine which computes the axis score deterministically:

\begin{mathbox}{Affinitree Score Formula}
\[
  S_\text{axis} \;=\; \sum_{i=1}^{k} w_i \cdot v_i \cdot m_i
\]
\begin{itemize}
  \item $w_i$ — concept weight (pre-calibrated constant, updated by ridge regression)
  \item $v_i$ — concept value from LLM scoring: $\hat{c}_i \in [0,1]$
  \item $m_i$ — evidence-tier multiplier: $0.6$ (T1, self-reported),
        $1.0$ (T2, artefact-backed), $1.2$ (T3, daemon-observed)
\end{itemize}
\end{mathbox}

\begin{itemize}
  \item The formula is a \emph{transparent linear combination}: every concept contributes
        a readable term $w_i \cdot v_i \cdot m_i$ whose value can be printed individually.
  \item The evidence-tier multiplier $m_i$ creates a direct incentive for founders to
        connect real tools. Linking GitHub makes a product concept T3, increasing its
        contribution by 100\% over the self-reported T1 baseline.
  \item The score is \textbf{fully deterministic}: given the same inputs the same result
        always follows. This allows Score Debate (human override vs.\ baseline) to be
        meaningful, since the baseline never fluctuates randomly.
\end{itemize}

\subsection{Bottleneck Identification}

The bottleneck is the concept with the smallest marginal contribution:
\[
  \text{bottleneck} \;=\; \arg\min_{i} \; (w_i \cdot v_i \cdot m_i)
\]

\begin{itemize}
  \item This identifies not the concept with the lowest absolute score, but the one
        that suppresses the axis total the most relative to its weight — the true weak link.
  \item Moufida surfaces this in the dashboard as: \emph{``Your weakest concept in [Axis]
        is [name]: contributing X\% of its potential. Improving it from 0.18 to 0.60
        would move your score from 2.3 to 3.1.''}
  \item The projected lift is computed by holding all other terms constant and solving
        for the new $S_\text{axis}$ at an improved $v_i$.
\end{itemize}

\subsection{Weight Calibration via Ridge Regression}

As the founder accumulates diagnostic history and a human reviewer locks score overrides,
the Signal service re-estimates concept weights to better reflect ground truth:

\begin{mathbox}{Ridge Calibration}
\[
  \hat{\mathbf{w}} \;=\; \arg\min_{\mathbf{w}}\;\left\| X\mathbf{w} - \mathbf{y} \right\|_2^2
  \;+\; \lambda\,\|\mathbf{w}\|_2^2
\]
\begin{itemize}
  \item $X \in \mathbb{R}^{n \times k}$ — matrix of concept scores across $n$ past diagnostics
  \item $\mathbf{y} \in \mathbb{R}^n$ — human-locked scores (expert override ground truth)
  \item $\lambda$ — ridge penalty that prevents weight collapse when data is sparse
\end{itemize}
\end{mathbox}

\begin{itemize}
  \item Ridge regression was chosen over plain least squares because early in a startup's
        history there are very few locked overrides ($n \ll k$). Ridge regularisation keeps
        weights well-conditioned and avoids fitting noise.
  \item After calibration, the Signal service updates weights in memory and persists them
        to disk. Future diagnostics automatically use the improved weights — a self-improving
        feedback loop tied directly to human expert judgement.
\end{itemize}

\subsection{How CBM Differs from LIME and SHAP}

\begin{center}\small
\begin{tabular}{@{}lll@{}}
\toprule
& \textbf{LIME / SHAP} & \textbf{Moufida CBM} \\
\midrule
Explanation unit        & Raw input features         & Human-named concepts \\
Causal structure        & Post-hoc approximation     & Intrinsic intermediate layer \\
Requires model internals & Yes (gradient/perturbation) & No \\
Adapts with feedback    & No                          & Yes (ridge recalibration) \\
Actionable output       & ``Field X matters''         & ``Improve concept Y by Z to gain $+P$ points'' \\
\bottomrule
\end{tabular}
\end{center}

\begin{implbox}{Implementation in Moufida}
\begin{itemize}[leftmargin=1em]
  \item Each axis service defines 4–7 concepts in natural language as part of its rubric schema.
  \item The orchestrator sends a structured concept-scoring prompt to \texttt{llama3.1:8b}
        with \texttt{temperature=0} and \texttt{seed=42} to ensure deterministic concept scores.
  \item The Affinitree engine (embedded in each axis service) receives the concept scores
        and evidence tiers, applies the formula, and returns the full breakdown alongside the score.
  \item The Signal service handles calibration and exposes \texttt{GET /cbm/weights/\{axis\}}
        and \texttt{POST /cbm/calibrate} to trigger a re-estimation pass when a new locked
        snapshot is recorded.
  \item The admin panel's LLM Calls view lets a judge verify that concept-scoring calls
        are real: one call per axis per diagnostic, with readable prompt preview and response.
\end{itemize}
\end{implbox}

%% ─────────────────────────────────────────────────────────────────────────────
\section{Layer 2 — Contrastive Axis Direction Probe}
%% ─────────────────────────────────────────────────────────────────────────────

\subsection{Academic Foundation}

\textbf{Zou et al.} — \textit{Representation Engineering: A Top-Down Approach to AI Transparency.}
arXiv 2023.\\
\textbf{Park et al.} — \textit{The Linear Representation Hypothesis and the Geometry of
Large Language Models.} arXiv 2023.

\begin{conceptbox}{The Core Idea}
Representation Engineering (RepE) proposes that high-level concepts are \emph{linearly
encoded} in the embedding space of a language model: there exists a unit direction vector
$\mathbf{d} \in \mathbb{R}^d$ such that documents about a concept project positively onto
$\mathbf{d}$ and unrelated documents project near zero.
Moufida applies this to the bge-m3 embedding space (1024-dim): one direction vector is
learned per diagnostic axis, encoding ``axis relevance'' as a geometric direction in the
shared multilingual embedding space.
\end{conceptbox}

\begin{notebox}
Standard RAG ranks chunks by cosine similarity to the query.
A document about ``startup funding rounds'' scores high for any query containing ``startup''
— even when the current axis is legal-compliance and the query is about IP.
Axis directions encode \emph{diagnostic relevance}, not topical overlap.
\end{notebox}

\subsection{Direction Extraction}

For each diagnostic axis, a set of contrastive document pairs is assembled:
$\{(x_j^+, x_j^-)\}_{j=1}^{N}$ where $x_j^+$ is clearly axis-relevant and $x_j^-$ is
a plausible but axis-irrelevant document (e.g., a legal text used as a negative example
for the market-intelligence axis).

The direction is extracted by contrasting the centroid embeddings:

\begin{mathbox}{Axis Direction Vector}
\[
  \boldsymbol{\mu}^+ = \frac{1}{N}\!\sum_j \phi(x_j^+),
  \qquad
  \boldsymbol{\mu}^- = \frac{1}{N}\!\sum_j \phi(x_j^-)
\]
\[
  \mathbf{d}_\text{axis} \;=\;
  \frac{\boldsymbol{\mu}^+ - \boldsymbol{\mu}^-}{\|\boldsymbol{\mu}^+ - \boldsymbol{\mu}^-\|_2}
  \;\in\; \mathbb{R}^{1024}
\]
where $\phi(\cdot)$ denotes the bge-m3 embedding function.
\end{mathbox}

\begin{itemize}
  \item The unit vector $\mathbf{d}_\text{axis}$ points in the direction of maximum contrast
        between axis-positive and axis-negative documents.
  \item Using centroid difference (mean of positives minus mean of negatives) makes the
        direction robust to individual outlier documents in the training corpus.
  \item Nine direction vectors are computed — one per diagnostic axis — and stored
        in memory inside the Signal service as 36 KB of floats at startup.
\end{itemize}

\subsection{Axis Relevance Scoring}

Given a candidate chunk embedding $\phi(c)$, the axis relevance score is:

\begin{mathbox}{Axis Relevance Score}
\[
  r_\text{axis}(c) \;=\;
  \sigma\!\!\left(\frac{\phi(c)\cdot\mathbf{d}_\text{axis} - b_\text{axis}}{\tau}\right)
  \;\in\; [0,1]
\]
where $b_\text{axis}$ is the mean dot product of neutral documents (calibrated bias),
$\sigma$ is the sigmoid function, and $\tau = 0.1$ is a temperature constant.
\end{mathbox}

\begin{itemize}
  \item The bias term $b_\text{axis}$ centres the score so that documents unrelated to any
        axis produce $r \approx 0.5$ rather than clustering at an arbitrary baseline.
  \item Temperature $\tau = 0.1$ was chosen to produce a sharp, discriminative distribution:
        axis-relevant chunks score $>0.8$ and irrelevant ones score $<0.3$.
\end{itemize}

\subsection{Hybrid RAG Re-ranking}

The RAG pipeline first retrieves top-$K$ candidates using Reciprocal Rank Fusion (RRF)
over dense (Qdrant) and sparse (BM25) rankings, then applies a final blend:

\begin{mathbox}{Final Re-ranking Score}
\[
  \text{score}_\text{final}(c) \;=\; 0.7 \cdot \text{RRF}(c) \;+\; 0.3 \cdot r_\text{axis}(c)
\]
Chunks matching the founder's sector tag receive an additional $1.3\times$ boost.
\end{mathbox}

\begin{itemize}
  \item The 70/30 split preserves retrieval diversity (RRF) while meaningfully shifting
        ranking toward diagnostically relevant documents.
  \item A pure axis-direction ranking would favour documents that are axis-relevant but
        not directly useful for the specific query — the blend prevents this regression.
  \item The sector boost applies as a multiplier on the final score, not as a hard filter,
        so non-sector documents remain retrievable when sector-specific evidence is absent.
\end{itemize}

\subsection{Two Uses in Moufida}

\begin{enumerate}
  \item \textbf{RAG re-ranking at diagnosis time.} Every axis service calls the RAG service
        with the current axis tag. The RAG service calls the Signal service
        (\texttt{POST /probe/project}) to score all candidate chunks against all 9 axes
        simultaneously, then blends the axis-specific score into the final ranking.

  \item \textbf{Auto-tagging at ingest time.} When a new document enters the knowledge
        base (via the admin panel or a file-drop from the companion), the top-2 axis
        relevance scores are stored as payload tags in Qdrant:
        \[
          \text{tags}(c) = \mathop{\mathrm{top}}_2\bigl\{r_\text{axis}(c) : \text{axis} \in A\bigr\}
        \]
        These tags act as fast pre-filters in subsequent retrievals.
\end{enumerate}

\begin{implbox}{Implementation in Moufida}
\begin{itemize}[leftmargin=1em]
  \item Direction vectors are computed offline by a one-time script using the curated KB
        as the positive corpus and cross-axis documents as negatives.
  \item They are stored on disk and loaded into the Signal service Rust process at startup.
        No re-computation occurs during a diagnostic — the probe lookup is a single dot
        product per chunk, O(1) in time and negligible in latency.
  \item bge-m3 was chosen because it is multilingual (Arabic, French, English) and
        produces 1024-dim embeddings with strong cross-lingual alignment. A Darija query
        and a French document co-exist in the same geometric space, so the direction vectors
        work across all three languages Moufida supports.
  \item The Signal service exposes \texttt{POST /probe/project} to score a batch of chunk
        embeddings for all 9 axes in a single call, minimising network round-trips.
\end{itemize}
\end{implbox}

%% ─────────────────────────────────────────────────────────────────────────────
\section{Layer 3 — Computerized Adaptive Testing (CAT / IRT)}
%% ─────────────────────────────────────────────────────────────────────────────

\subsection{Academic Foundation}

\textbf{Lord} — \textit{Applications of Item Response Theory to Practical Testing Problems.} 1977.\\
\textbf{Bock \& Mislevy} — \textit{Adaptive EAP Estimation of Ability in a Microcomputer
Environment.} Applied Psychological Measurement, 1982.\\
\textbf{Wainer et al.} — \textit{Computerized Adaptive Testing: A Primer.} 2nd ed., 2000.

\begin{conceptbox}{The Core Idea}
Traditional questionnaires ask every founder the same fixed set of questions.
Item Response Theory (IRT) observes that different questions are informative at different
ability levels: a question about cap-table structure is nearly irrelevant for an
ideation-stage founder, while a question about whether the founder has a co-founder
tells an advanced team nothing new.
Computerized Adaptive Testing (CAT) selects each next question to \emph{maximally
reduce uncertainty} about the founder's latent startup maturity $\theta$, producing a
more precise estimate in fewer questions.
\end{conceptbox}

\subsection{Two-Parameter Logistic Model (2PL IRT)}

The probability that a founder with latent maturity $\theta$ gives a positive answer to
item $j$ follows a logistic curve:

\begin{mathbox}{2PL Item Response Function}
\[
  P(u_j = 1 \mid \theta) \;=\; \frac{1}{1 + e^{-a_j(\theta - b_j)}}
\]
\begin{itemize}
  \item $\theta \in \mathbb{R}$ — founder's latent startup maturity (prior mean at $0$)
  \item $a_j > 0$ — discrimination: how sharply item $j$ separates founders around $b_j$
  \item $b_j \in \mathbb{R}$ — difficulty: $\theta$ at which the founder has a 50\% chance
        of a positive answer
\end{itemize}
\end{mathbox}

\begin{itemize}
  \item The 2PL model was chosen over the simpler 1PL (Rasch) model because items vary
        significantly in how sharply they discriminate: a question about revenue
        ($a_j \approx 2.1$) separates founders much more cleanly than one about business
        name registration ($a_j \approx 0.8$).
  \item Items with high discrimination and matched difficulty provide the most information
        about the current ability estimate — which is exactly what CAT exploits.
\end{itemize}

\subsection{Item Selection via Fisher Information}

\begin{mathbox}{Fisher Information and Selection Rule}
\[
  I_j(\theta) \;=\; a_j^2 \cdot P_j(\theta) \cdot \bigl[1 - P_j(\theta)\bigr]
\]
\[
  j^* \;=\; \arg\max_{j \,\notin\, \text{administered}}\; I_j(\hat{\theta}_n)
\]
\end{mathbox}

\begin{itemize}
  \item Fisher information peaks when $P_j(\theta) = 0.5$, i.e., when the item difficulty
        matches the current ability estimate. This is why CAT never asks questions that are
        too easy or too hard: they add little information.
  \item Maximising Fisher information at each step is equivalent to minimising expected
        posterior variance, making the procedure Bayesian-consistent.
\end{itemize}

\subsection{Bayesian EAP Estimation}

After each response, the ability estimate is updated using the Expected A Posteriori rule:

\begin{mathbox}{EAP Update}
\[
  \hat{\theta}_{n+1} \;=\;
  \frac{%
    \displaystyle\int \theta \cdot L(\mathbf{u}_n \mid \theta) \cdot p(\theta)\, d\theta
  }{%
    \displaystyle\int L(\mathbf{u}_n \mid \theta) \cdot p(\theta)\, d\theta
  }
\]
\[
  L(\mathbf{u}_n \mid \theta) \;=\; \prod_{j=1}^n P_j(\theta)^{u_j}\,[1-P_j(\theta)]^{1-u_j}
\]
Prior: $p(\theta) \sim \mathcal{N}(0,1)$.
Integral approximated on a 61-point quadrature grid $\theta \in [-4,\,4]$.
\end{mathbox}

\begin{itemize}
  \item The standard normal prior represents the assumption that Tunisian founders are
        broadly distributed around an average maturity level — neither systematically
        advanced nor systematically nascent.
  \item Bayesian EAP was chosen over Maximum Likelihood Estimation because MLE is undefined
        when all early responses are consistently positive or negative, which happens
        frequently with ideation-stage founders. EAP handles this gracefully via the prior.
  \item The 61-point quadrature provides sufficient numerical precision for the 0–12 item
        range used in practice while remaining fast enough to recompute on every request.
\end{itemize}

\subsection{Stopping Rule and Two-Phase Item Bank}

The session stops when either:
\begin{enumerate}
  \item $\text{SE}(\hat{\theta}) < 0.40$ — posterior uncertainty is sufficiently low, \emph{or}
  \item $n \geq 12$ — hard cap to prevent session fatigue.
\end{enumerate}

This yields \textbf{8–15 questions} in practice, vs.\ 30+ for a fixed survey
at equivalent measurement precision.

\begin{center}\small
\begin{tabular}{@{}lll@{}}
\toprule
\textbf{Phase} & \textbf{Items} & \textbf{Purpose} \\
\midrule
Stage discriminators & 15 binary items & Coarse placement: ideation / MVP / growth / scale \\
Axis-specific items  & 8–12 per axis  & Fine-grained targeting given the stage estimate \\
\bottomrule
\end{tabular}
\end{center}

\begin{implbox}{Implementation in Moufida}
\begin{itemize}[leftmargin=1em]
  \item The CAT engine lives entirely inside the orchestrator and is \textbf{stateless}:
        the client sends the full response history $(j_1,u_1),\dots,(j_n,u_n)$ on every
        call to \texttt{POST /intake/next}. There is no server-side session. This was a
        deliberate architectural choice for horizontal scalability — any orchestrator
        replica can serve the same session without shared state.
  \item The item bank is a static JSON file loaded into memory at startup. Each item
        carries its $a_j$ and $b_j$ parameters, phase tag, and axis affiliation.
  \item The orchestrator maps the final $\hat{\theta}$ to a discrete maturity stage using
        fixed thresholds, then converts the answer history into a profile patch that
        pre-fills structured fields before the founder continues to the dashboard.
  \item The intake session is interpretable in a strict sense: every question selection
        is provably optimal given the response history. An auditor can replay the selection
        algorithm with the same history and verify that the same question would have been chosen.
\end{itemize}
\end{implbox}

%% ─────────────────────────────────────────────────────────────────────────────
\section{How the Three Layers Reinforce Each Other}
%% ─────────────────────────────────────────────────────────────────────────────

The three interpretability mechanisms operate at different stages and compound:

\begin{center}\small
\begin{tabular}{@{}llll@{}}
\toprule
\textbf{Layer} & \textbf{Pipeline stage} & \textbf{Output} & \textbf{Consumed by} \\
\midrule
CAT/IRT  & Intake (pre-diagnosis) & Richer profile patch, $\hat{\theta}$ & CBM concept scoring \\
Probe    & RAG retrieval          & Axis-relevant chunks                 & All axis LLM prompts \\
CBM      & Post-scoring           & Bottleneck + calibrated score        & Dashboard UI \\
\bottomrule
\end{tabular}
\end{center}

\begin{itemize}
  \item A better intake (CAT) produces more informative profile fields, which yields
        more coherent LLM concept scores in the CBM pass and more precise axis queries
        for the Probe to re-rank.
  \item Better retrieved evidence (Probe) improves the quality of LLM justifications,
        making the concept scores in the CBM pass more grounded in real Tunisian resources.
  \item Calibrated weights (CBM) reflect the Tunisian ecosystem context over time —
        concepts that genuinely predict success in Tunisia earn higher weights, creating
        a virtuous cycle as more diagnostics accumulate.
\end{itemize}

The combination is what produces full interpretability: the founder can trace from the
questions asked, through the evidence retrieved, through the concept breakdown, all the
way to the final score and projected improvement path — at every step with a computable,
auditable justification.

\end{document}
